\section{The Cognitive Manifold}
\label{sec:cognitive-manifold}

NDH interaction occurs on a geometric space we term the \emph{cognitive
manifold}, denoted $\mathcal{M}_{cog}$. This manifold is distinct from the
runtime manifold $M$ defined in the NDH core geometry. Whereas $M$ encodes the
internal modal structure and recursion kernel of NDH, the cognitive manifold
$\mathcal{M}_{cog}$ encodes the geometry of interrogation: how semantic inputs
are represented, how traversal unfolds, and how spectral components respond to
excitation.

\subsection{Definition}

The cognitive manifold is a smooth, connected manifold equipped with a
connection $\nabla^{cog}$ governing interaction transport:
\[
\mathcal{M}_{cog} \quad \text{smooth, connected}, \qquad
\nabla^{cog} : T\mathcal{M}_{cog} \to T\mathcal{M}_{cog}.
\]
Points $x \in \mathcal{M}_{cog}$ correspond to interrogation states of NDH.
Semantic inputs induce tangent vectors at these points, initiating geodesic
interaction trajectories.

\subsection{Modal Adjacency}

The modal partition of NDH,
\[
M = M_A \cup M_B \cup M_C,
\]
induces a corresponding partition of $\mathcal{M}_{cog}$:
\[
\mathcal{M}_{cog} = \mathcal{M}_A \cup \mathcal{M}_B \cup \mathcal{M}_C,
\]
where each region $\mathcal{M}_i$ represents the interaction domain associated
with modal region $M_i$. Modal boundaries in $M$ induce hypersurfaces in
$\mathcal{M}_{cog}$ across which interaction trajectories may undergo spectral
or recursion-dependent deformation.

\subsection{Tangent Structure}

Semantic excitation is modeled as a tangent vector
\[
\mathbf{s} \in T_x \mathcal{M}_{cog},
\]
representing the direction and magnitude of interrogation. The tangent bundle
$T\mathcal{M}_{cog}$ therefore encodes the local interaction geometry, and the
connection $\nabla^{cog}$ determines how excitation propagates along geodesics.

\subsection{Geodesic Initiation}

Given an excitation vector $\mathbf{s}$, the induced geodesic is defined by
\[
\gamma(t) = \exp_x(t\mathbf{s}),
\]
where $\exp_x$ is the exponential map associated with $\nabla^{cog}$. The
trajectory $\gamma(t)$ describes the evolution of NDH's interaction state under
semantic input. Modal-boundary crossings along $\gamma(t)$ correspond to
structural transitions in spectral emission and recursion alignment.

\subsection{Relation to Runtime Geometry}

The cognitive manifold is not identical to the runtime manifold $M$, but the
two are coupled. Runtime geometry determines the modal structure and recursion
kernel, while cognitive geometry determines how interrogation traverses these
structures. The connection $\nabla^{cog}$ is constrained by the holonomy clamp
defined on $M$, ensuring that interaction transport respects NDH's internal
geometry.

This coupling yields a coherent geometric model of NDH interaction, enabling
formal analysis of semantic excitation, geodesic traversal, and spectral bundle
behavior.
